\phantomsection
\chapter*{Kavramlar Sözlüğü}
\markboth{Kavramlar Sözlüğü}{Kavramlar Sözlüğü}
\addcontentsline{toc}{chapter}{Kavramlar Sözlüğü}

Bu sözlük, kitabın temel kavramlarını alfabetik olarak ve ilgili bölüm
bağlamıyla birlikte açıklar. Sayfa numarası yerine bölüm adı kullanılmıştır;
böylece PDF düzeni değişse bile kavramın ana tartışma yeri korunur. Kavramların
metin içinde geçtiği sayfalar kitabın sonundaki Teknik Dizin'de verilmiştir.

\begin{multicols}{2}
\raggedcolumns
\raggedright
\footnotesize
\setlength{\parindent}{0pt}
\setlength{\parskip}{0.22em}

\noindent\textbf{\large A}\par\smallskip
\textbf{A/B testi} --- İki alternatif kararın veya oran farkının veriyle karşılaştırılması; Bayesçi güncelleme ve klasik test bağlamında kullanılır. \emph{(4, 9)}\par
\textbf{Aile bazında hata oranı} --- Çoklu karşılaştırmalarda en az bir yanlış red olasılığını kontrol eden hata ölçüsüdür. \emph{(4, 6)}\par
\textbf{Akaike/Bilgi ölçütü} --- Model karşılaştırmada uyum ve karmaşıklık dengesini düşünmeye yarayan ölçüt ailesidir. \emph{(5, 9)}\par
\textbf{Aleatorik belirsizlik} --- Sistemin doğal rastgeleliğinden gelen ve veri artsa bile tamamen kaybolmayan belirsizliktir. \emph{(1)}\par
\textbf{Alternatif hipotez} --- Sıfır hipotezine karşı sınanan araştırma iddiasıdır. \emph{(4)}\par
\textbf{Anakütle} --- Çıkarım yapılmak istenen hedef birimlerin tümüdür. \emph{(1)}\par
\textbf{Anlamlılık düzeyi} --- Testte I. tip hata olasılığı için önceden belirlenen $\alpha$ sınırıdır. \emph{(4)}\par
\textbf{ANOVA} --- Üç veya daha fazla grup ortalamasını grup içi değişkenlikle karşılaştıran yöntemdir. \emph{(6)}\par
\textbf{Artık} --- Gözlenen değer ile modelin tahmin ettiği değer arasındaki farktır. \emph{(5)}\par
\textbf{Artık analizi} --- Model varsayımlarını ve uyum sorunlarını grafik veya sayısal tanılarla inceleme sürecidir. \emph{(5)}\par
\textbf{Aykırı değer} --- Veri deseninden belirgin biçimde ayrılan gözlemdir; otomatik silinmez, nedeni araştırılır. \emph{(2, 5, 8)}\par
\medskip
\noindent\textbf{\large B}\par\smallskip
\textbf{Bağımsız gözlemler} --- Bir gözlemin diğerinin değerini sistematik olarak belirlemediği varsayımdır. \emph{(1, 4, 7, 8)}\par
\textbf{Bayes tahmin edicisi} --- Sonsal beklenen kaybı küçükleyen tahmindir; kare hata kaybında sonsal ortalamadır. \emph{(9)}\par
\textbf{Bayesçi çıkarım} --- Önsel bilgi ile olabilirliği sonsal dağılımda birleştiren çıkarım yaklaşımıdır. \emph{(9)}\par
\textbf{Beklenen değer} --- Rassal değişkenin uzun dönemli merkezini ifade eden ortalama ölçüsüdür. \emph{(Ek A)}\par
\textbf{Beklenen frekans} --- Ki-kare testinde model veya bağımsızlık varsayımı altında beklenen hücre sayısıdır. \emph{(7)}\par
\textbf{Beta-Binom modeli} --- Binom verisi ile Beta önselinin eşlenik güncelleme modelidir. \emph{(9)}\par
\textbf{Bootstrap} --- Örneklemden yeniden örnekleme yaparak belirsizlik ve güven aralığı tahmin etme yöntemidir. \emph{(3)}\par
\medskip
\noindent\textbf{\large C--Ç}\par\smallskip
\textbf{Cohen'in $d$ değeri} --- Ortalama farkını standart sapma ölçeğinde veren etki büyüklüğüdür. \emph{(4)}\par
\textbf{Çarpanlara ayırma teoremi} --- Yeterli istatistiği ortak yoğunluğu $g(T,\theta)h(x)$ biçiminde ayırarak belirler. \emph{(1)}\par
\textbf{Çoklu karşılaştırma} --- ANOVA veya çoklu test sonrası hangi farkların anlamlı olduğunu kontrollü biçimde inceleme sürecidir. \emph{(4, 6)}\par
\medskip
\noindent\textbf{\large D}\par\smallskip
\textbf{Dağılım tablosu} --- Normal, $t$, ki-kare ve $F$ kritik değerlerine hızlı erişim sağlayan başvuru tablosudur. \emph{(Ek B)}\par
\textbf{Değişen varyans} --- Hata varyansının gözlem veya tahmin düzeyiyle birlikte değişmesi durumudur. \emph{(5)}\par
\textbf{Duyarlılık analizi} --- Sonuçların varsayım, önsel veya model seçimine ne kadar bağlı olduğunu inceleme sürecidir. \emph{(3, 9)}\par
\medskip
\noindent\textbf{\large E}\par\smallskip
\textbf{EKK} --- En küçük kareler yöntemi; regresyon katsayılarını artık kareler toplamını küçülterek tahmin eder. \emph{(5)}\par
\textbf{Eşlenik önsel} --- Önsel ve olabilirlik birleştiğinde aynı dağılım ailesinden sonsal üreten önseldir. \emph{(9)}\par
\textbf{Etki büyüklüğü} --- Bulguların pratik veya bilimsel önemini nicelleştiren ölçüdür. \emph{(4, 6, 7, 8)}\par
\textbf{Etkin örneklem hacmi} --- Karmaşık tasarım veya ağırlıklandırma sonrası bilgi içeriğini eşdeğer bağımsız örneklem sayısı olarak yorumlar. \emph{(1, 3)}\par
\textbf{Epistemik belirsizlik} --- Bilgi eksikliği, sınırlı örneklem veya model basitleştirmesinden kaynaklanan belirsizliktir. \emph{(1)}\par
\medskip
\noindent\textbf{\large F}\par\smallskip
\textbf{$F$ dağılımı} --- Varyans oranı, ANOVA ve regresyon testlerinde kullanılan dağılımdır. \emph{(6, Ek B)}\par
\textbf{Fisher-Neyman teoremi} --- Yeterlilik için çarpanlara ayırma ölçütünü veren temel teoremdir. \emph{(1)}\par
\medskip
\noindent\textbf{\large G}\par\smallskip
\textbf{Güç} --- Alternatif doğruyken $H_0$ hipotezini reddetme olasılığıdır. \emph{(4)}\par
\textbf{Güven aralığı} --- Frekansçı uzun dönem kapsama oranı üzerinden yorumlanan belirsizlik aralığıdır. \emph{(3)}\par
\textbf{Güvenilir aralık} --- Bayesçi sonsal olasılık yorumu taşıyan parametre aralığıdır. \emph{(9)}\par
\medskip
\noindent\textbf{\large H}\par\smallskip
\textbf{Hellinger uzaklığı} --- Dağılımlar arası farklılığı ölçen sağlam tahmin bağlamında kullanılan uzaklık ölçüsüdür. \emph{(2)}\par
\textbf{Hipotez testi} --- Verinin sıfır hipoteziyle uyumunu test istatistiği ve p-değeriyle değerlendiren çıkarım yöntemidir. \emph{(4)}\par
\textbf{HPD aralığı} --- Sonsal yoğunluğu en yüksek bölgeleri kapsayan Bayesçi güvenilir aralıktır. \emph{(9)}\par
\medskip
\noindent\textbf{\large I--İ}\par\smallskip
\textbf{I. tip hata} --- Doğru $H_0$ hipotezini reddetme hatasıdır. \emph{(4)}\par
\textbf{II. tip hata} --- Yanlış $H_0$ hipotezini reddedememe hatasıdır. \emph{(4)}\par
\textbf{İstatistik} --- Örneklemden hesaplanan ve parametre içermeyen fonksiyondur. \emph{(1)}\par
\textbf{İstatistiksel çıkarım} --- Örneklemden anakütle veya model hakkında belirsizlik içeren sonuç üretme sürecidir. \emph{(1)}\par
\medskip
\noindent\textbf{\large K}\par\smallskip
\textbf{Karar kuralı} --- Test, tahmin veya model seçimi sonucunda hangi eylemin seçileceğini belirleyen kuraldır. \emph{(4, 9)}\par
\textbf{Ki-kare testi} --- Gözlenen ve beklenen frekanslar arasındaki farkı değerlendiren kategorik veri testidir. \emph{(7)}\par
\textbf{Kayıp fonksiyonu} --- Karar veya tahmin hatasının maliyetini ölçen fonksiyondur. \emph{(9)}\par
\textbf{Olumsallık tablosu} --- İki veya daha fazla kategorik değişkenin ortak frekans tablosudur. \emph{(7)}\par
\textbf{Korelasyon} --- İki değişken arasındaki doğrusal ilişkinin yön ve gücünü ölçer. \emph{(5, Ek A)}\par
\textbf{Kovaryans} --- İki değişkenin birlikte değişimini ölçen momenttir. \emph{(Ek A)}\par
\textbf{Kritik değer} --- Red bölgesini belirleyen dağılım kantilidir. \emph{(4, Ek B)}\par
\medskip
\noindent\textbf{\large M}\par\smallskip
\textbf{Mann--Whitney testi} --- İki bağımsız grubun sıra temelli konum farkını inceleyen parametrik olmayan testtir. \emph{(8)}\par
\textbf{MAP tahmini} --- Sonsal yoğunluğu en büyükleyen Bayesçi nokta tahminidir. \emph{(9)}\par
\textbf{MSE} --- Ortalama karesel hata; varyans ve yanlılığı birlikte içeren tahmin kalitesi ölçüsüdür. \emph{(2)}\par
\textbf{Momentler yöntemi} --- Dağılım momentlerini örneklem momentlerine eşitleyerek parametre tahmin eden yöntemdir. \emph{(2)}\par
\medskip
\noindent\textbf{\large N}\par\smallskip
\textbf{Neyman--Pearson yaklaşımı} --- Hipotez testini hata olasılıkları ve güç üzerinden kuran klasik çerçevedir. \emph{(4)}\par
\textbf{Normal dağılım} --- Güven aralığı, test ve regresyon hataları için merkezi rol oynayan sürekli dağılımdır. \emph{(3, 4, 5, Ek B)}\par
\medskip
\noindent\textbf{\large O--Ö}\par\smallskip
\textbf{Olabilirlik} --- Verinin parametre değerlerine göre göreli desteğini gösteren fonksiyondur. \emph{(2, 9)}\par
\textbf{Önsel dağılım} --- Bayesçi analizde veri öncesi parametre belirsizliğini ifade eder. \emph{(9)}\par
\textbf{Örneklem} --- Anakütleden gözlenen veya seçilen birimler kümesidir. \emph{(1)}\par
\textbf{Örnekleme dağılımı} --- Bir tahmin edicinin tekrarlı örnekleme altında sahip olduğu dağılımdır. \emph{(2, 3)}\par
\textbf{Örnekleme tasarımı} --- Anakütleden birim seçme planıdır; standart hata ve genelleme gücünü etkiler. \emph{(1, 3)}\par
\medskip
\noindent\textbf{\large P}\par\smallskip
\textbf{p-değeri} --- $H_0$ doğruyken gözlenen kadar veya daha aşırı sonuç elde etme olasılığıdır. \emph{(4)}\par
\textbf{Parametre} --- Anakütle veya modelin bilinmeyen sabit özelliğidir. \emph{(1, 2)}\par
\textbf{Parametrik olmayan test} --- Dağılım varsayımı zayıf olduğunda kullanılan sıra veya işaret temelli test ailesidir. \emph{(8)}\par
\textbf{Pearson katkısı} --- Ki-kare istatistiğinde bir hücrenin toplam istatistiğe katkısıdır. \emph{(7)}\par
\medskip
\noindent\textbf{\large R}\par\smallskip
\textbf{Rao--Blackwell iyileştirmesi} --- Bir tahmin ediciyi yeterli istatistiğe göre koşullu bekleyerek varyansını azaltma yöntemidir. \emph{(2)}\par
\textbf{Raporlama kalıbı} --- Bulguların tahmin, belirsizlik, varsayım ve yorumla birlikte yazılmasını sağlayan standart ifade yapısıdır. \emph{(Tüm bölümler)}\par
\textbf{Regresyon} --- Bağımlı değişkenin açıklayıcı değişkenlerle ortalama ilişkisinin modellenmesidir. \emph{(5)}\par
\textbf{Risk} --- Karar kuramında kayıp fonksiyonunun beklenen değeridir. \emph{(2, 9)}\par
\medskip
\noindent\textbf{\large S--Ş}\par\smallskip
\textbf{Serbestlik derecesi} --- Test istatistiğinin dağılımını belirleyen bağımsız bilgi miktarıdır. \emph{(4, 6, 7, Ek B)}\par
\textbf{Sıfır hipotezi} --- Testte başlangıçta doğru kabul edilerek sınanan temel iddiadır. \emph{(4)}\par
\textbf{Sıra istatistiği} --- Örneklem gözlemlerinin küçükten büyüğe sıralanmış değerleridir. \emph{(1, 8)}\par
\textbf{Sonlu anakütle düzeltmesi} --- Anakütlenin önemli kısmı örneklendiğinde standart hatayı azaltan düzeltmedir. \emph{(1, 3)}\par
\textbf{Sonsal dağılım} --- Bayesçi analizde önsel ile olabilirliğin birleşiminden elde edilen veri sonrası dağılımdır. \emph{(9)}\par
\textbf{SPSS} --- Menü tabanlı uygulama ve çıktı yorumlama için kullanılan istatistik yazılımıdır. \emph{(Tüm bölümler)}\par
\textbf{Standart hata} --- Tahmin edicinin örnekleme dağılımındaki standart sapmadır. \emph{(2, 3)}\par
\textbf{Student $t$ dağılımı} --- Varyans bilinmediğinde ortalama çıkarımı için kullanılan dağılımdır. \emph{(3, Ek B)}\par
\textbf{Şekil/grafik yorumu} --- Formül veya karar kuralının görsel davranışını inceleme pratiğidir. \emph{(Tüm bölümler)}\par
\medskip
\noindent\textbf{\large T}\par\smallskip
\textbf{Tahmin} --- Gözlenen veriden hesaplanan sayısal parametre değeridir. \emph{(2)}\par
\textbf{Tahmin edici} --- Örneklemden parametre tahmini üreten rassal kuraldır. \emph{(2)}\par
\textbf{Tasarım etkisi} --- Karmaşık örnekleme tasarımının varyansı basit rastgele örneklemeye göre kaç kat değiştirdiğini gösterir. \emph{(1, 3)}\par
\textbf{Tek yönlü ANOVA} --- Tek faktörlü grup karşılaştırmalarında kullanılan ANOVA modelidir. \emph{(6)}\par
\textbf{Tutarlılık} --- Örneklem hacmi arttıkça tahmin edicinin hedef parametreye yaklaşmasıdır. \emph{(2)}\par
\medskip
\noindent\textbf{\large U--Ü}\par\smallskip
\textbf{Uygunluk testi} --- Gözlenen frekansların belirli bir dağılım veya oran yapısıyla uyumunu sınar. \emph{(7)}\par
\textbf{Üstel dağılım} --- Bekleme zamanı ve güvenilirlik problemlerinde kullanılan sürekli dağılımdır. \emph{(1, 2)}\par
\medskip
\noindent\textbf{\large V}\par\smallskip
\textbf{Varyans} --- Bir rassal değişkenin beklenen değeri etrafındaki değişkenliğini ölçer. \emph{(Ek A)}\par
\textbf{Varyans homojenliği} --- Grupların ortak hata varyansına sahip olması varsayımıdır. \emph{(6)}\par
\textbf{Varsayım kontrolü} --- Yöntemin geçerli yorumlanması için gereken koşulları analiz öncesi ve sonrası değerlendirme sürecidir. \emph{(Tüm bölümler)}\par
\medskip
\noindent\textbf{\large W--Y--Z}\par\smallskip
\textbf{Welch ANOVA} --- Varyans homojenliği zayıfladığında grup ortalamalarını karşılaştırmak için kullanılan sağlam ANOVA alternatifidir. \emph{(6)}\par
\textbf{Wilcoxon testi} --- Eşleşmiş veya sıra tabanlı konum karşılaştırmalarında kullanılan parametrik olmayan testtir. \emph{(8)}\par
\textbf{Yanlılık} --- Tahmin edicinin beklenen değeri ile hedef parametre arasındaki sistematik farktır. \emph{(2)}\par
\textbf{Yansızlık} --- Tahmin edicinin beklenen değerinin hedef parametreye eşit olmasıdır. \emph{(2)}\par
\textbf{Yeterli istatistik} --- Parametre hakkında örneklemdeki bilgiyi koruyan veri indirgeme istatistiğidir. \emph{(1)}\par
\textbf{Yöntem seçim kontrol listesi} --- Araştırma sorusu, varsayım, veri tipi, belirsizlik ve raporlama adımlarını birlikte denetleyen pratik listedir. \emph{(Tüm bölümler)}\par
\textbf{Z testi} --- Varyans bilindiğinde veya büyük örneklemde kullanılan normal yaklaşım testidir. \emph{(4, Ek B)}\par

\normalsize
\end{multicols}